% !TEX root = ../main.tex
\section{Three Teams, Three Budgets}
\label{sec:three-teams}

Everything above is stated as though an organisation were a single agent that could
decide to satisfy Table~\ref{tab:invariants-system}. It is not, and the way it is
divided is itself a principal cause of the failures this paper describes. This
section is the least formal in the paper and, for practitioners, possibly the most
useful.

\subsection{The division}

Responsibility for the conditions of warrant is distributed across organisational
functions that were not designed to hold it jointly. In financial services the fourth
row below is a second-line function with an independent reporting line and a
supervisory mandate; in other sectors the equivalent may sit in quality assurance,
clinical governance, or internal audit. Wherever it sits, it is the function that
performs the independent challenge this paper argues for, and it is the one most often
absent from discussions of these systems.

\begin{table}[ht]
  \centering
  \small
  \begin{tabular}{@{}>{\raggedright\arraybackslash}p{2.8cm}>{\raggedright\arraybackslash}p{4.6cm}>{\raggedright\arraybackslash}p{2.4cm}>{\raggedright\arraybackslash}p{4.0cm}@{}}
    \toprule
    Function & Typical responsibility & Owns & Failure if left in isolation \\
    \midrule
    Data / knowledge
      & Source ingestion, curation, lineage, metadata, quality
      & I8, I9
      & Authentic documents with no policy applicability or output control \\
    \addlinespace
    Policy / risk / compliance
      & Rules, jurisdiction, decision authority, approvals, procedures
      & I10, I15
      & Governing policy with no control over how assertions are retrieved,
        generated, or verified \\
    \addlinespace
    Platform / AI engineering
      & Models, prompts, orchestration, retrieval, evaluation, monitoring
      & I11--I14, I17
      & Answer quality, latency, and usability with no evidentiary authority or
        institutional authorization \\
    \addlinespace
    Model risk / independent validation
      & Independent challenge, validation, and approval of models and their use
      & I13, I15, I16
      & Validating $M$ while the risk sits in $S$: a well-documented model, an
        unexamined evidence path \\
    \bottomrule
  \end{tabular}
  \caption{Four functions, and the invariants each is positioned to supply. No
    function owns a complete path from source to authorized reliance.}
  \label{tab:three-teams}
\end{table}

These are not merely three technical workstreams. They typically have separate
budgets, separate incentives, separate release cadences, separate accountability
structures, and separate systems of record. A data team is measured on lineage
coverage and freshness. A compliance function is measured on policy completeness
and audit findings. A platform team is measured on answer quality, latency,
adoption, and cost. None is measured on the grounding rate of
\eqref{eq:rate}, because it spans all three.

\subsection{The unit of failure is the handoff}

The consequence is that failures concentrate at the boundaries rather than inside
any function's remit. A system may have:

\begin{itemize}
  \item excellent data lineage and weak policy binding---I8 and I9 without I10;
  \item well-specified governing policy and uncontrolled retrieval---I10 without
        I11;
  \item high benchmark accuracy and no per-assertion provenance---nothing in
        Table~\ref{tab:invariants-system} at all, with a validation report that
        looks complete;
  \item citations without applicability validation---I12 without I11, which is
        precisely cell (v);
  \item human review without a reproducible evidence path---review that cannot be
        re-performed and therefore cannot be audited, which is I13 without I14;
  \item a validated model inside an unvalidated system---documentation of $M$'s
        training, evaluation, and limitations, with no examination of context
        resolution, applicability logic, defeater checking, or authorization. This is
        the characteristic failure of the fourth row, and it is the most dangerous
        because it produces a formal approval.
\end{itemize}

Each of these is a system in which every function has done its job. The ungoverned
handoff is the operational unit of failure, and it is invisible from inside any one
function's metrics.

This also explains a pattern worth naming, because it is the most expensive form of
the problem. Cell (v)---misapplied grounding---sits exactly on the boundary between
the data function, which supplies authentic versioned sources, and the policy
function, which knows which sources govern which decisions. The metadata required to
enforce applicability is a policy artefact stored in a data system and consumed by a
platform component. It is nobody's deliverable, and it is the single most common
thing I find missing.

\subsection{What follows}

Two conclusions, and the first is uncomfortable for how these programmes are
usually organised.

Evidentiary controls must be architecture and cross-functional invariants, not
after-the-fact review tasks. A review step owned by one function cannot establish a
property that requires artefacts from two others, and adding reviewers does not
change this. Table~\ref{tab:invariants-system} is a specification for a system, and
a specification that spans three budgets needs an owner who spans them.

A fifth function appears once retention is considered. Privacy and records
management own the tension described in Section~\ref{sec:retention}: the audit record
that reconstructability requires is in direct conflict with minimisation and erasure
obligations, and neither the data, policy, platform, nor validation function is
positioned to resolve it. That handoff has the same structure as the others and is
usually discovered later.

And the grounding rate is the right cross-functional metric precisely because no
single function can move it alone. Accuracy can be improved by the platform team
without reference to the others, which is part of why it is the number that gets
reported. \eqref{eq:rate} cannot: improving it requires corpus work, policy
metadata, and pipeline control together. A metric that only improves when three
functions cooperate is an unusually honest instrument for a problem whose failures
live between them.
